\documentclass[10pt,a4paper]{article}

\usepackage[utf8]{inputenc}
\usepackage[T1]{fontenc}
\usepackage[english]{babel}   % or english
\usepackage{hyperref}

\usepackage{minted}
\usepackage{xcolor}

% Nice inline code / commit hashes
\newcommand{\commit}[1]{\texttt{\textcolor{gray}{#1}}}

\title{Slipshow: chill-coding with OCaml}
\author{Paul-Elliot}
\date{\today}

\begin{document}
\maketitle

\section{Introduction}

Slipshow is a tool to make presentation. It tries to innovate from the
traditional ``deck of slides'' that rules the current presentations. For instance, it introduces the concept of slips, that are slides without a bottom limit, on
  which you scroll to display the content. Another feature that is unique to Slipshow, as far as I know, is the ability to record a drawing, edit it, and replay it during the
  presentation.
% \item A mix between the WYSIWYG and WYSIWYM model, where the latter one is used
%   for most style and layout, and the former one for anything that has to be
%   placed absolutely.

\section{Growing the project}

Slipshow started as a single file containing some JavaScript, that could be
included in an HTML page written in a specific way, to turn it into a
presentation.  However, this very quickly became unmaintainable. Basic features
were buggy and endlessly fixed, breaking other part of the program.

It is now a relatively big project, that features a runtime, a compiler, an LSP
server with a preview window, a record-and-replay for hand drawings with a
specialized user interface to edit them. How did we go from a simpler thing
being unmaintainable, to the current version?

\subsection{Chill-coding}

Today, the answer for such question is very often ``vibe-coding''. Vibe-coding
is the practice of relying entirely on AI for generating the code, without even
reading it. Freed from the friction of interacting with the code, considered too
low-level, the cognitive load is lowered. You can start building with your vibe.

At first, from the name ``vibe-code'' and the good vibes that I get while
coding, I thought I was practicing it. The first issue is that I am not using AI
to generate any code. The second issue is that vibe-coding implies losing the
knowledge of the internals of your software, which can make it hard to really
trust enough to distribute it, at least not without really extensive testing.

I like to know the internals of my softwares, and more than vibes, I feel
\emph{chill} while coding. Chill when doing big refactors. Chill when adding new
features. Chill for feeling safe from introducing bugs or complexity to the
codebase. Which is why I introduce the term \emph{chill-coding}, for the act of
coding without fear.

\emph{Chill-coding} requires either exceptional skills, or support from powerful
tools. I fall into the second category.  After being stuck for years with the
JavaScript prototype, I moved to OCaml as the primary language for
Slipshow. Relying on the excellent compiler errors, state-of-the-art editor
support, exceptional library ecosystem and over-the-top community, I was able to
\emph{chill-code} all the way.

\section{An experience report}

I have been very happy with the experience of writing in OCaml since the
beginning. As mentioned above, it feels safe and lets your creativity wander,
without friction. My hope is to share a bit of this excitement for the language,
the compiler and the tools around, as well as the OCaml ecosystem and community,
through a selected (git) history of Slipshow.

\subsection{A runtime in JavaScript}

As mentioned, Slipshow started as a presentation engine written in
JavaScript. Here is an excerpt of commits all focused on fixing the same basic
feature: going back in the presentation:

\commit{3af56a3} --- \texttt{corrected horribug previous}\\
\commit{436bb65} --- \texttt{implemented engine.previousSlip ← the feature is born}\\
\commit{d0cfc69} --- \texttt{corrected "not coming beginning on previous" bug}\\
\commit{61e6cb7} --- \texttt{fix "future slip, delay not 0 again when previousing"}\\
\commit{3af56a3} --- \texttt{corrected horribug previous}\\
\commit{bcd8323} --- \texttt{better previous: use delay from previous state}\\
\commit{28372cd} --- \texttt{two bugs in previous and annotations}\\
\commit{2437301} --- \texttt{correcting bug with "previous" and "resizeObserver"}

% The single-JS-file era: a bug swamp.

\subsection{OCaml for the compiler}

In \commit{0fa489e}, the first line of OCaml is introduced in the project. It is
the commit introducing the compiler, that takes a markdown file and outputs a
standalone HTML file.

OCaml is known to be very good for compiler implementation. It is no exception
in this case, especially when relying on \texttt{cmdliner} for the command line,
\texttt{grace} for error reporting, and \texttt{cmarkit} for Markdown parsing.

One relatively exotic feature from OCaml that I am benefiting from are
extensible variants. \texttt{cmarkit} defines its AST using extensible variants,
allowing the user of the library to extend it. Note however that with extensible
variants, one loses some kind chill points, as exhaustiveness checks cannot be
made for extensible variants. In commit \commit{5713489} could get the chill back with:

\begin{minted}{ocaml}
type s_block =
| Slip of Block.t attributed node
| SlipScript of Block.Code_block.t attributed node
| [...]

type Cmarkit.Block.t += S_block of s_block
\end{minted}

\commit{5713489} --- \texttt{Make it more suitable for match exhaustiveness check}.

% OCaml enters: Markdown $\to$ slip show.

\subsection{The runtime, rewritten in OCaml}

For some time, the JavaScript code for the presentation engine, and the OCaml code for the compiler, co-existed. However, I still had the issues mentioned above on maintainability, and making the two communicate was challenging. So it was time to rewrite it all in OCaml.

The best way to summarize the rewriting is the following commit:

\commit{32d0917} --- \emph{``BOOOOOOOOOOOOOOOMM!'' (Undo monad, they work)}

This commit is about the same feature that the commits mentioned in section~3.1 were trying to fix. But I've already spoke about the Undo monad enough, so in this section, I'll focus on how \texttt{dune}, \texttt{js\_of\_ocaml} and
\texttt{ppx\_blob} made it absolutely trivial to drive the build of Slipshow,
which is more complex than usual: The Slipshow runtime is compiled to bytecode,
then to JavaScript, and finally embedded in the source of the Slipshow compiler.

\begin{minted}{lisp}
  (executable
    (name engine)
    (modes js))

  (executable
    (name compiler)
    (preprocess (pps ppx_blob))  ; Use with [%blob "engine.bc.js"]
    (preprocessor_deps engine.bc.js))
\end{minted}


\subsection{Record and replay drawings}

One recent and game-changing addition to Slipshow is the ability to record
drawings, edit them and replay the result during the presentation.

The UI to edit records, make very much like a video editor, despite the expected challenge to craft, proved to be a joy to
make. This joy can be summed up with the following commit:

\commit{3ac99e3} --- \emph{``Wo so cool lwd''}

Again, the joy comes from high quality libraries from the ecosystem. Lwd is a library for reactive programming. It allows to define the UI from reactive values, and have the UI update as the reactive values change. Here, we shorten path to what was written at \texttt{elpased\_time}:

\begin{minted}{ocaml}
let filter_path elapsed_time path =
  let$* path = path in  (* the two inputs are reactive values *)
  let$ elapsed_time = elapsed_time in
  List.filter (fun (_, t) -> t < elapsed_time) path
\end{minted}

% \subsection{An LSP server}

% The youngest addition to Slipshow is an LSP server. It felt extremely easy
% thanks to \texttt{linol}, a library from the ecosystem. It is nice that it is
% using objects, adding them to my collections of OCaml features.

% \commit{e568193} --- \emph{``LSP: actually openings of new files already work''}

\section{Conclusion}

If Slipshow was able to grow, if I am able to chill-code it and have fun all days, this is thanks to an incredible language, compiler, ecosystem and community. I hope to be able to send this message and share my excitement.

\end{document}
